\section{Polygon size and the role of a marked edge}\label{sec:introduction}

A large self-avoiding polygon can have many more edges than its diameter.
The question is how these two sizes are related when each polygon receives
its critical weight.  We answer it for the regular honeycomb lattice by
estimating three positive sums: the mass of large polygons, the mass seen
from one marked edge, and the mass seen from two marked edges.

Use the unit-edge lattice with vertices
\[
 P_{ij}=ih+jv,\qquad R_{ij}=P_{ij}+(h+v)/3,\qquad
 h=(\sqrt3,0),\quad v=(\sqrt3/2,3/2),\quad i,j\in\mathbb Z.
\]
The neighbours of $P_{ij}$ are $R_{ij},R_{i-1,j},R_{i,j-1}$.
A polygon is a finite simple unoriented cycle.  We identify polygons that
differ by a translation in $\mathbb Zh+\mathbb Zv$, and count each resulting
class once.  For such a class $\gamma$, write $L(\gamma)$ for its number of
edges and $D(\gamma)$ for its Euclidean diameter.  Its critical weight is
\[
 w(\gamma)=x_*^{L(\gamma)},\qquad x_*=(2+\sqrt2)^{-1/2}.
\]
The value of $x_*$ is the inverse honeycomb connective constant proved by
Duminil-Copin and Smirnov~\cite{DuminilCopinSmirnov2012}.

\begin{theorem}[Diameter tails and the second length moment]\label{thm:polygon-tails}
\label{thm:polygon-second-moment}
For the polygon classes and weights just defined, put
\[
 A(r)=\sum_{D(\gamma)\ge r}w(\gamma),\qquad
 P(r)=\sum_{D(\gamma)\ge r}L(\gamma)w(\gamma),\qquad
 M_2(r)=\sum_{D(\gamma)\le r}L(\gamma)^2w(\gamma).
\]
Then, as $r\to\infty$,
\begin{equation}\label{eq:main-polygon}
 A(r)=r^{-2+o(1)},\qquad P(r)=r^{-2/3+o(1)},\qquad
 M_2(r)\le r^{2/3+o(1)}.
\end{equation}
Every sum in this statement is finite.  Here $f(r)=r^{a+o(1)}$ means
that for every $\varepsilon>0$, $r^{a-\varepsilon}\le f(r)\le
r^{a+\varepsilon}$ for all sufficiently large $r$; the one-sided notation
has the analogous meaning.
\end{theorem}

The ratio $P(r)/A(r)$ is the mean length under the unrooted critical law
conditioned on diameter at least $r$.  Thus this mean has exponent $4/3$.
The second moment in \eqref{eq:main-polygon} gives additional information:
it prevents too much critical mass from being carried by very long
polygons of small diameter.  Combining it with the two tails proves the
following laws for positive cutoff ensembles.

\begin{corollary}[Length cutoffs]\label{cor:length-cutoffs}
The length tails satisfy
\[
 \sum_{L(\gamma)\ge n}w(\gamma)=n^{-3/2+o(1)},\qquad
 \sum_{L(\gamma)\ge n}L(\gamma)w(\gamma)=n^{-1/2+o(1)}.
\]
Under either normalized weight $w(\gamma)$ or $L(\gamma)w(\gamma)$,
conditioned on $L(\gamma)\ge n$, one has
$D(\gamma)=n^{3/4+o(1)}$ in probability.  The same conclusion holds
with the length restriction $n\le L(\gamma)<n^{1+\delta}$ for any
fixed $\delta>0$.  Under the length-weighted law conditioned on
$D(\gamma)\ge r$, one has $L(\gamma)=r^{4/3+o(1)}$ in probability.
An assertion $X_t=t^{a+o(1)}$ in probability means that its value lies
between $t^{a-\varepsilon}$ and $t^{a+\varepsilon}$ with probability
tending to one for every $\varepsilon>0$.
\end{corollary}

The choice of weight matters: the mean statement for unrooted polygons
conditioned on large diameter and the typical-length statement for
length-weighted polygons have different normalizations.

There are also open-path consequences.  Deleting a prescribed edge from a
polygon yields a walk between its two endpoints; the resulting fixed-edge
law is a constant multiple of the length-weighted polygon law for
symmetry-invariant restrictions. A bounded modification at a lowest
vertex transfers the unrooted estimates to half-plane walks between two
prescribed boundary vertices at fixed separation. Their diameter-conditioned
mean is stated in Corollary~\ref{cor:boundary-neighbor-mean};
Section~\ref{sec:open-cutoffs} gives the length-cutoff consequences.

The final part of the paper proves a separate family of boundary
estimates. Its paths have endpoints at edge midpoints, called ports,
and their weights count visited vertices, so concatenation at a port
multiplies weights exactly. Theorem~\ref{thm:boundary-arches} shows that
the critical mass of half-plane arches between ports at separation $n$
is comparable to $n^{-5/4}$, with a comparable mass confined to diameter
$O(n)$. Theorem~\ref{thm:slab-corridors} shows that the crossing mass
from a fixed lower port to a free upper port of a strip of height $H$
is comparable to $H^{-1/4}$. The same order of mass remains inside any
prescribed piecewise linear guide corridor of fixed positive relative
width. The precise domains and confinement conditions are given with
those theorems.

A further confinement consequence concerns the convergence radius $x_{\rm strip}(W)$ of walks restricted to a vertical strip of half-width $W$:
\[
 0\le x_{\rm strip}(W)-x_*\le W^{-4/3+o(1)}.
\]
Proposition~\ref{prop:strip-radius} proves this one-sided estimate by stacking polygons with recoverable joins.

\subsection{History and the main difficulty}

Nienhuis's analysis of the dilute $O(n)$ model predicted the critical
point and planar self-avoiding-walk exponents, including the size
exponent $3/4$~\cite{Nienhuis1982}. Lawler, Schramm and Werner later
formulated scaling-limit conjectures for walks and polygons and recovered
this exponent from the conjectural connection with
$\mathrm{SLE}_{8/3}$~\cite{LawlerSchrammWerner2004}.
Corollary~\ref{cor:length-cutoffs} establishes this size relation for the
critical polygon cutoff ensembles defined above. The passage from
positive tail sums to typical size requires no scaling-limit assumption;
conditioning on one exact length requires additional local control.

Duminil-Copin and Smirnov established the honeycomb connective constant
using a parafermionic observable whose local identity has a positive
boundary consequence~\cite{DuminilCopinSmirnov2012}. For boundary paths,
Beaton, Bousquet-M\'elou, de Gier, Duminil-Copin and Guttmann proved that
the critical mass of bridges crossing a strip tends to zero with its
height~\cite{BeatonEtAl2014}. Krachun and Panagiotis proved a polynomial
upper bound for this mass and an $O(n/\log n)$ radius bound with high
probability for the uniform $n$-step honeycomb
walk~\cite{KrachunPanagiotis2026}. Their geometric estimate concerns a
fixed-length walk law. Our boundary results give matching constant-factor
bounds with the predicted bridge exponent $-1/4$, including confinement
to prescribed corridors.

The local Yang--Baxter weights used here belong to the dilute-loop
integrable setting. The bulk analysis of Zhou and Batchelor and the
rhombic self-avoiding-walk work of Glazman and Manolescu provide the
relevant local-weight and transport
framework~\cite{ZhouBatchelor1997,GlazmanManolescu2017}.
Garbali and Nienhuis study polynomial exchange equations and fusion
recurrences for dilute $O(1)$ ground states, with determinant formulas
for associated sum rules~\cite{GarbaliNienhuis2017Ground,GarbaliNienhuis2017Sum}.
We construct the two zero-loop-weight vectors needed here by
interpolation and identify their physical normalization. The angular
construction uses the classical lattice vertex-operator
method~\cite{FrenkelKac1980,Segal1981,Frenkel1985}. The bilinear form in
our construction is indefinite, so we first work coefficientwise before
establishing analytic continuation. The radial estimates use a separate
positive Hilbert-space realization.

The main analytic obstacle is that the useful transformed expressions
are complex.  Their terms are neither probabilities nor positive
weights.  A formal leading exponent therefore does not determine the
physical polygon sum: its coefficient might vanish.  We resolve this by
combining exact identities with positive geometric information obtained
before the transformed asymptotic analysis.  The compact operators used
later act on positive Hilbert spaces, but their kernels remain complex;
their spectral information comes from exact trace identities.

\subsection{From finite cylinders to positive polygons}

Project the lattice to a cylinder of circumference $N|h|$.  A planar
polygon either remains simple after projection or suffers a collision
with a translate.  The first class corresponds to contractible cylinder
polygons.  Cylinder polygons winding around the circumference supply a
second positive observable.  The proof compares these two classes using
finite transfer matrices before taking any limit.

The cylinder calculation also gives positive leading constants: the winding mass per class modulo translation along the cylinder axis is asymptotic to $c/N$ with $c>0$, and the unrooted mass lost under planar projection is asymptotic to a positive constant times $N^{-2}$. Its transverse-span estimate, measuring extent along the cylinder axis, permits arbitrary polynomial suppression, with any fixed length weight, for both winding and contractible polygons (Proposition~\ref{prop:thin-cylinder-span}). These estimates retain information beyond the tail exponents.

The finite calculation uses two polynomial fixed vectors, with seam
twists $\eta=-i$ and $\eta=i$. At physical weights their marked
contractions have a simple interpretation: the first counts winding
plus contractible polygons through the marked edge, and the second
counts winding minus contractible polygons. Adding them therefore
isolates winding mass. Let $E_N$ be the critical winding mass through
three sampled edges, one of each honeycomb edge type, counted once for
each sampled edge used. Proposition~\ref{a:finite-polygons} gives
\[
                 E_N\tau_N=(2-\sqrt2)\tau^v_{N-1}.
\]
Here $\tau_N$ and $\tau^v_{N-1}$ are the homogeneous continued values of
two finite color sums. Their individual terms need not be positive.
The determinant calculation below proves that $\tau_N$ is nonzero, so
the identity expresses the positive mass $E_N$ as their ratio.

Sections~\ref{c4:angular}--\ref{c4:plateau} estimate that ratio. The
angular expansion records the algebra and phases of labeled
insertions; the radial form places the insertions at real positions.
A radial configuration also carries a piecewise constant slope in a
specified lattice. The slope changes by a prescribed amount at each
insertion, and its squared size gives an exponential cost along the
radial coordinate. The central marked charge suppresses certain
insertion types over a region of logarithmic width. A coercive estimate
controls the absolute sum despite the complex phases and gives compact
transfer operators for consecutive intervals. The required jump
between the allowed slopes on the two sides determines the marked
decay relative to the unmarked determinant. The positive winding bound
from Section~\ref{c4:pressure} prevents the leading coefficient from
vanishing. The ratio has exponent $-2/3$, and
Section~\ref{c4:tails} converts it and the unrooted projection loss into
the two planar diameter tails.

For two marks, let $k-1$ be the shorter site block between the marked
columns and let $N$ be the total number of sites. Under the hypotheses
of Theorem~\ref{c4:thm-two-mark}, the central comparison region has
width of order $\log k$. When the two block sizes differ, an
intermediate region remains in which only the larger charge sector is
suppressed. Sections~\ref{c4:two-mark}--\ref{c4:hierarchical} derive the
two-mark finite identity and Pfaffian normalization, prove coercivity
for the doubled charge system, and bound the trace in that intermediate
region. The three contributions combine as
\[
       k^{-9/8}(N/k)^{5/24}N^{-5/24}=k^{-4/3}.
\]
Section~\ref{c4:metric-consequences} realizes the permitted arrays by
tilted cylinders, identifies the scalar with the positive mass through
two prescribed edges, and sums over their displacement. Short
displacements and three lattice directions are handled directly by the
first length moment. This proves the truncated second moment.
Section~\ref{sec:cutoff-laws} then derives the cutoff laws using only
the three positive polygon estimates.

The boundary arguments use the finite vacuum with twist $\eta=-i$ and
its even Pfaffian normalization. A three-root identity controls the mass
of pairs of arches that intersect. Combined with the cylinder arch loss,
it gives the flat-boundary pointwise and spatial bounds. The slab
construction then retains positive mass in prescribed guide corridors,
controlling the multiplicity of intermediate cuts before removing those
marks.

Throughout, constants in power bounds may depend on a fixed exponent
slack and on fixed geometric parameters.  Algebraic notation is local
to its section; in particular the auxiliary bilinear spaces introduced
for the angular calculation are distinct from the physical plane.
